\documentclass[11pt,a4paper]{article}
\usepackage[utf8]{inputenc}
\usepackage[T1]{fontenc}
\usepackage[french]{babel}
\usepackage{geometry}
\geometry{margin=2.5cm}
\usepackage{graphicx}
\usepackage{xcolor}
\usepackage{amsmath}
\usepackage{listings}
\usepackage{hyperref}
\usepackage{fancyhdr}
\usepackage{titlesec}
\usepackage{enumitem}
\usepackage{tcolorbox}
\usepackage{float}
\usepackage{tikz}
\usepackage{tabularx}
\usepackage{longtable}
\usepackage{booktabs}
\usepackage{fontawesome5}
\usetikzlibrary{shapes.geometric,arrows.meta,positioning,calc,fit,backgrounds}
\usepackage{amssymb}
\usepackage[table]{xcolor}   % pour gérer les couleurs dans les tableaux
\usepackage{colortbl}        % pour activer \rowcolor



% Configuration des listes
\setlist[itemize]{leftmargin=*, itemsep=3pt, topsep=3pt}
\setlist[enumerate]{leftmargin=*, itemsep=3pt, topsep=3pt}

% Configuration du code
\lstset{
    basicstyle=\ttfamily\small,
    breaklines=true,
    frame=single,
    tabsize=4,
    showstringspaces=false,
    captionpos=b
}

% En-tête et pied de page
\pagestyle{fancy}
\fancyhf{}
\fancyhead[L]{\small Roadmap Prototype Chatbot SUP'PTIC}
\fancyhead[R]{\small 8-13 février 2026}
\fancyfoot[C]{\thepage}

% Environnements personnalisés
\newtcolorbox{important}[1][]{
  colback=red!5!white,
  colframe=red!75!black,
  fonttitle=\bfseries,
  title=#1
}

\newtcolorbox{info}[1][]{
  colback=blue!5!white,
  colframe=blue!75!black,
  fonttitle=\bfseries,
  title=#1
}

\newtcolorbox{astuce}[1][]{
  colback=green!5!white,
  colframe=green!75!black,
  fonttitle=\bfseries,
  title=#1
}

\newtcolorbox{exemple}[1][]{
  colback=orange!5!white,
  colframe=orange!75!black,
  fonttitle=\bfseries,
  title=#1
}

\begin{document}

% Page de titre
\begin{titlepage}
    \centering
    \vspace*{2cm}
    {\Huge \textbf{ROADMAP PROTOTYPE}} \\[0.5cm]
    {\LARGE \textbf{CHATBOT SUP'PTIC}} 
    
    \vspace{1cm}
    
    {\Large 8 février - 13 février 2026 (6 jours)}
    
    \vspace{1.5cm}
    
    \begin{minipage}{0.8\textwidth}
        \centering
        \small Documentation technique complète pour \\
        le développement du prototype
    \end{minipage}
    
    \vfill
    
    \begin{center}
        \textit{Rédigé par :} \\
        \vspace{0.3cm}
        \textbf{NKOUMOU TJADE} \\
        Chef de la Division de Programmation \\
        Club Informatique SUP'PTIC
    \end{center}
    
\end{titlepage}

% Table des matières
\tableofcontents
\clearpage

% Section 1 : Contexte
\section{Contexte et Contraintes}
\label{sec:contexte}

\subsection{Situation actuelle}

\begin{itemize}
    \item[\faIcon{check-circle}] \textbf{Documentation technique complète} : algorithme TF-IDF, architecture BD, gestion FAQ v2.0
    \item[\faIcon{check-circle}] \textbf{Maquette wireframe} déjà proposée (à finaliser)
    \item[\faIcon{times-circle}] \textbf{Pas de données CSV réelles} : nécessité de génération massive
\end{itemize}

\subsection{Équipes disponibles (10 personnes)}

\begin{enumerate}
    \item \textbf{Base de Données (2 places)} \\
    Mission : Django models, migrations, optimisation, scripts d'import
    
    \item \textbf{Structuration des Données (4 places)} \\
    Mission : Génération massive CSV avec IA, nettoyage, validation
    
    \item \textbf{Backend (2 places)} \\
    Mission : API Django, TF-IDF, similarité cosinus, endpoints REST
    
    \item \textbf{Frontend (2 places)} \\
    Mission : Interface chat HTML/CSS/JS, connexion API, design responsive
\end{enumerate}

\subsection{Stratégie de Génération des Données}

\begin{important}[Approche par Vagues de 100 Q/R]
L'équipe Structuration des Données fonctionne en \textbf{vagues de 100 Q/R par sous-thème}.

\textbf{Objectif ambitieux :} 1000+ Q/R pour le prototype

\textbf{Méthode :}
\begin{itemize}
    \item Chaque sous-groupe génère 100 Q/R par session de travail
    \item Focus sur un sous-thème à la fois pour garantir cohérence
    \item Utilisation intensive d'IA générative (ChatGPT, Claude, Copilot)
    \item Validation par échantillonnage (20\% des Q/R)
\end{itemize}
\end{important}

\clearpage

% Section 2 : Objectif
\section{Objectif du Prototype}
\label{sec:objectif}

\subsection{Livrable final (13 février 18h)}

Un chatbot fonctionnel démontrant :

\begin{itemize}
    \item[\faIcon{check}] \textbf{Recherche par similarité TF-IDF} opérationnelle sur 1000+ Q/R
    \item[\faIcon{check}] \textbf{Base de données Django} avec modèles complets et relations
    \item[\faIcon{check}] \textbf{API REST fonctionnelle} : endpoints FAQ, feedback, statistiques
    \item[\faIcon{check}] \textbf{Interface HTML/CSS/JS} basée sur wireframe existant
    \item[\faIcon{check}] \textbf{1000+ Q/R validées} couvrant toutes les catégories principales
    \item[\faIcon{check}] \textbf{Système de feedback} opérationnel (\faThumbsUp \, \faThumbsDown)
    \item[\faIcon{check}] \textbf{Documentation complète} : README, API, guide BD
\end{itemize}

\subsection{Principe de Génération Massive}

\begin{info}[Pipeline de Production]
\textbf{Jour 1-2 :} 400 Q/R \\
\textbf{Jour 3-4 :} 600 Q/R supplémentaires \\
\textbf{Total :} 1000+ Q/R en 4 jours

\vspace{0.3cm}
\textbf{Répartition par catégorie :}
\begin{itemize}
    \item Introduction \& Historique : 100 Q/R
    \item Inscriptions \& Admissions : 100 Q/R
    \item Examens \& Évaluations : 100 Q/R
    \item Vie Étudiante : 100 Q/R
    \item Services Administratifs : 100 Q/R
    \item Programmes Académiques : 100 Q/R
    \item Infrastructures \& Ressources : 100 Q/R
    \item Carrières \& Insertion Professionnelle : 100 Q/R
    \item Communication \& Vie Institutionnelle : 100 Q/R
    \item Vie Pratique : 100 Q/R
\end{itemize}

\textbf{Total :} 1000 Q/R
\end{info}

\clearpage

% Section 3 : Planning
\section{Planning Détaillé Jour par Jour}
\label{sec:planning}

\subsection{DIMANCHE 8 FÉVRIER - Jour 1 : Fondations et Première Vague}

\subsubsection{Matin (8h-12h) - Setup \& Démarrage}

\paragraph{Équipe Base de Données}
\begin{itemize}
    \item[$\square$] Créer nouveau dépôt Git propre
    \item[$\square$] Initialiser projet Django (django-admin startproject)
    \item[$\square$] Définir les 5 models Django :
    \begin{itemize}
        \item Utilisateurs (extension AbstractUser)
        \item Catégories
        \item FAQ (question, réponse, catégorie, sous-thème, source)
        \item FAQVector (vecteur TF-IDF, norme)
        \item Feedback (utilisateur, FAQ, type, score, commentaire)
    \end{itemize}
    \item[$\square$] Ajouter contraintes et relations (ForeignKey, CASCADE/RESTRICT)
    \item[$\square$] Première migration (makemigrations + migrate)
\end{itemize}

\paragraph{Équipe Structuration des Données}
\begin{itemize}
    \item[$\square$] \textbf{Sous-groupe 1 :} Première vague de génération
    \begin{itemize}
        \item ≥ 100 Q/R sur "Introduction \& Historique"
        \item Utiliser le CSV de référence déjà défini
        \item Respecter les règles de qualité et validation
    \end{itemize}

    \item[$\square$] \textbf{Sous-groupe 2 :} Deuxième vague de génération
    \begin{itemize}
        \item ≥ 100 Q/R sur "Inscriptions \& Admissions"
        \item Respecter le même format CSV de référence
        \item Validation croisée entre sous-groupes
    \end{itemize}
\end{itemize}


\paragraph{Équipe Backend}
\begin{itemize}
    \item[$\square$] Installer dépendances (requirements.txt)
    \begin{lstlisting}[language=bash]
Django==4.2.7
djangorestframework==3.14.0
django-cors-headers==4.3.1
scikit-learn==1.3.2
pandas==2.1.3
numpy==1.26.2
    \end{lstlisting}
    \item[$\square$] Créer structure apps Django : \texttt{faq/}, \texttt{chatbot/}, \texttt{users/}
    \item[$\square$] Implémenter module prétraitement texte (tokenization, stopwords français)
\end{itemize}

\paragraph{Équipe Frontend}
\begin{itemize}
    \item[$\square$] Finaliser wireframe de l'interface chat
    \item[$\square$] Setup projet HTML/CSS/JS :
    \begin{itemize}
        \item Structure de base : index.html, styles.css, app.js
        \item Framework CSS (Bootstrap)
    \end{itemize}
    \item[$\square$] Créer composants de base (zone messages, input, boutons)
\end{itemize}

\subsubsection{Après-midi (14h-18h) - Production Intensive}

\paragraph{Tous ensemble}
\begin{itemize}
    \item[$\square$] \textbf{Réunion synchronisation} (14h-14h30)
    \item[$\square$] Officialiser architecture projet sur nouveau dépôt
    \item[$\square$] Commit initial avec structure complète
\end{itemize}

\textbf{Architecture projet}
\begin{verbatim}
ChatBot/
├── backend/
│   ├── config/          # Paramètres Django (settings, urls, wsgi)
│   ├── faq/             # App gestion FAQ (stockage et réponses)
│   ├── chatbot/         # App algorithme TF-IDF et prétraitement texte
│   ├── users/           # App gestion des utilisateurs + feedback
│   ├── manage.py        # Commandes Django (runserver, migrations)
│   └── db.sqlite3       # Base de données SQLite par défaut
│
├── frontend/            # Interface utilisateur (HTML/CSS/JS)
│   ├── index.html
│   ├── css/             # Styles
│   ├── js/              # Scripts dynamiques
│   └── assets/          # Images, icônes, ressources statiques
│
├── data/
│   ├── csv/             # Fichiers CSV générés par structuration
│   └── scripts/         # Scripts import/export (CSV ↔ BD)
│
├── docs/                # Documentation technique et fonctionnelle
│
├── venv/                # Environnement virtuel Python (isolé)
│
└── requirements.txt     # Liste des dépendances Python

\end{verbatim}


\paragraph{Équipe Structuration (Production massive)}
\begin{itemize}
    \item[$\square$] \textbf{Sous-groupe 1 :} Troisième vague
    \begin{itemize}
        \item ≥ 100 Q/R sur "Examens \& Évaluations - Modalités"
        \item Respecter le format CSV de référence
        \item Validation interne avant livraison
    \end{itemize}

    \item[$\square$] \textbf{Sous-groupe 2 :} Quatrième vague
    \begin{itemize}
        \item ≥ 100 Q/R sur "Vie Étudiante - Clubs et associations"
        \item Respecter le même format CSV de référence
        \item Validation interne avant livraison
    \end{itemize}

    \item[$\square$] \textbf{Validation croisée}
    \begin{itemize}
        \item Chaque sous-groupe valide 20\% des Q/R de l'autre
        \item Vérification cohérence, clarté et absence de doublons
    \end{itemize}
\end{itemize}


\paragraph{Équipe Backend}
\begin{itemize}
    \item[$\square$] Implémenter fonction \texttt{preprocess\_text(text)}
    \item[$\square$] Créer fixtures de test (10 FAQ, 3 catégories, 2 utilisateurs)
\end{itemize}

\paragraph{Équipe Frontend}
\begin{itemize}
    \item[$\square$] Implémenter structure HTML complète
    \item[$\square$] Styliser interface (CSS)
    \item[$\square$] Préparer logique d'envoi requête (fetch API)
\end{itemize}

\paragraph{Équipe Base de Données}
\begin{itemize}
    \item[$\square$] Créer script d'import CSV → Django
    \item[$\square$] Tester import avec 50 premières Q/R
    \item[$\square$] Créer index sur colonnes critiques
\end{itemize}

\subsubsection{Fin de journée (18h-18h30)}
\begin{itemize}
    \item[$\square$] \textbf{Stand-up :} chaque équipe présente avancement
    \item[$\square$] Vérifier objectif atteint : 300 Q/R générées
    \item[$\square$] Identifier blocages éventuels
\end{itemize}

\clearpage

\subsection{LUNDI 9 FÉVRIER - Jour 2 : Algorithme et Vague Massive}

\subsubsection{Matin (8h-12h) - Cœur Algorithmique}

\paragraph{Équipe Backend}
\begin{itemize}
    \item[$\square$] Implémenter vectorisation TF-IDF complète :
    \begin{itemize}
        \item Fonction \texttt{train\_vectorizer(corpus)}
        \item Fonction \texttt{compute\_tfidf\_vector(text)}
        \item Calcul et stockage vecteurs dans FAQVector
        \item Calcul normes vectorielles
    \end{itemize}
    \item[$\square$] Implémenter similarité cosinus :
    \begin{itemize}
        \item Fonction \texttt{compute\_cosine\_similarity(vec1, vec2)}
        \item Fonction \texttt{find\_best\_faq(question, top\_k=3)}
    \end{itemize}
\end{itemize}

\paragraph{Équipe Structuration (Vagues 4-5-6)}
\begin{itemize}
    \item[$\square$] \textbf{Vague 5 :} 100 Q/R "Services Administratifs - Scolarité"
    \begin{itemize}
        \item Respecter le format CSV de référence
        \item Validation interne avant livraison
    \end{itemize}

    \item[$\square$] \textbf{Vague 6 :} 100 Q/R "Programmes Académiques - Filières"
    \begin{itemize}
        \item Respecter le même format CSV de référence
        \item Validation interne avant livraison
    \end{itemize}

    \item[$\square$] \textbf{Objectif session :} 300 Q/R supplémentaires → Total : 600 Q/R
    \begin{itemize}
        \item Vérifier cohérence et homogénéité des données
        \item Validation croisée entre sous-groupes
    \end{itemize}
\end{itemize}


\paragraph{Équipe Base de Données}
\begin{itemize}
    \item[$\square$] Optimiser script d'import (traitement par batch)
    \item[$\square$] Import des 600 Q/R dans la base
    \item[$\square$] Vérifier intégrité référentielle
\end{itemize}

\paragraph{Équipe Frontend}
\begin{itemize}
    \item[$\square$] Implémenter logique envoi requête
    \item[$\square$] Affichage réponse chatbot (bulle de message)
    \item[$\square$] Gérer états de chargement (spinner)
\end{itemize}

\subsubsection{Après-midi (14h-18h) - Intégration \& Production}

\paragraph{Équipe Backend}
\begin{itemize}
    \item[$\square$] Créer endpoints API REST (Django REST Framework) :
    \begin{itemize}
        \item \texttt{POST /api/chatbot/ask/} - poser question
        \item \texttt{GET /api/faq/} - lister FAQ
        \item \texttt{GET /api/categories/} - lister catégories
    \end{itemize}
    \item[$\square$] Tester endpoints avec Postman/curl
\end{itemize}

\paragraph{Équipe Structuration (Vagues 7-8)}
\begin{itemize}
    \item[$\square$] \textbf{Vague 7 :} 100 Q/R "Carrières \& Insertion - Stages"
    \begin{itemize}
        \item Respecter le format CSV de référence
        \item Validation interne avant livraison
    \end{itemize}

    \item[$\square$] \textbf{Vague 8 :} 100 Q/R "Inscriptions - Documents requis"
    \begin{itemize}
        \item Respecter le même format CSV de référence
        \item Validation interne avant livraison
    \end{itemize}

    \item[$\square$] \textbf{Validation finale :} vérifier cohérence 800 Q/R totales
    \begin{itemize}
        \item Contrôle qualité global
        \item Vérification absence de doublons
        \item Validation croisée entre sous-groupes
    \end{itemize}
\end{itemize}


\paragraph{Équipe Frontend}
\begin{itemize}
    \item[$\square$] Connecter frontend à API backend
    \item[$\square$] Tester envoi question et affichage réponse
    \item[$\square$] Gérer erreurs (pas de réponse, API indisponible)
\end{itemize}

\paragraph{Équipe Base de Données}
\begin{itemize}
    \item[$\square$] Import des 800 Q/R totales
    \item[$\square$] Vérifier que tous vecteurs sont calculés
    \item[$\square$] Requêtes de vérification (count, duplicates)
\end{itemize}

\subsubsection{Soir (18h-19h)}
\begin{itemize}
    \item[$\square$] \textbf{Test end-to-end} : parcours utilisateur complet
    \item[$\square$] Vérifier : 800 Q/R en base, algorithme fonctionnel
    \item[$\square$] Identifier bugs critiques à corriger demain
\end{itemize}

\clearpage

\subsection{MARDI 10 FÉVRIER - Jour 3 : Feedback et Finalisation 1000+}

\subsubsection{Matin (8h-12h) - Système Feedback}

\paragraph{Équipe Backend}
\begin{itemize}
    \item[$\square$] Implémenter endpoint feedback :
    \begin{itemize}
        \item \texttt{POST /api/feedback/} - enregistrer feedback
        \item Sauvegarder : user\_id, faq\_id, type (\faThumbsUp/\faThumbsDown), score, commentaire
    \end{itemize}
    \item[$\square$] Créer endpoint statistiques :
    \begin{itemize}
        \item \texttt{GET /api/stats/} - FAQ par taux satisfaction
        \item \texttt{GET /api/stats/categories/} - répartition par catégorie
    \end{itemize}
\end{itemize}

\paragraph{Équipe Structuration (Vagues finales 9-10)}
\begin{itemize}
    \item[$\square$] \textbf{Vague 9 :} 100 Q/R "Examens - Rattrapages et notes"
    \begin{itemize}
        \item Respecter le format CSV de référence
        \item Validation interne avant livraison
    \end{itemize}

    \item[$\square$] \textbf{Vague 10 :} 100 Q/R "Vie Pratique - Transport, restauration"
    \begin{itemize}
        \item Respecter le même format CSV de référence
        \item Validation interne avant livraison
    \end{itemize}

    \item[$\square$] \textbf{OBJECTIF ATTEINT :} 1000 Q/R
    \begin{itemize}
        \item Vérification cohérence globale
        \item Contrôle qualité final (échantillon 10\%)
        \item Enrichissement : ajouter variantes pour Q/R importantes
    \end{itemize}
\end{itemize}


\paragraph{Équipe Base de Données}
\begin{itemize}
    \item[$\square$] Ajouter 5-10 utilisateurs de test
    \item[$\square$] Créer données feedback fictives pour démo
    \item[$\square$] Tester requêtes analytiques (statistiques)
\end{itemize}

\paragraph{Équipe Frontend}
\begin{itemize}
    \item[$\square$] Ajouter boutons feedback (\faThumbsUp/\faThumbsDown) sous réponses
    \item[$\square$] Implémenter boîte commentaire feedback négatif
    \item[$\square$] Améliorer design (CSS responsive)
\end{itemize}

\subsubsection{Après-midi (14h-18h) - Raffinement}

\paragraph{Équipe Backend}
\begin{itemize}
    \item[$\square$] Ajouter système seuil de confiance :
    \begin{itemize}
        \item Score < 0.6 : "Je n'ai pas trouvé de réponse précise"
        \item 0.6 < Score < 0.8 : "Voici la réponse la plus proche"
        \item Score > 0.8 : Réponse directe
    \end{itemize}
    \item[$\square$] Optimiser recherche (caching si nécessaire)
\end{itemize}

\paragraph{Équipe Structuration}
\begin{itemize}
    \item[$\square$] Documenter toutes sources utilisées (SOURCES.md)
    \item[$\square$] Créer CSV final propre et consolidé
    \item[$\square$] Validation qualité finale (échantillon 10\%)
\end{itemize}

\paragraph{Équipe Frontend}
\begin{itemize}
    \item[$\square$] Ajouter page "Statistiques" basique
    \item[$\square$] Implémenter liste catégories cliquables
    \item[$\square$] Améliorer UX mobile (responsive design)
\end{itemize}

\paragraph{Équipe Base de Données}
\begin{itemize}
    \item[$\square$] Import final 1000+ Q/R
    \item[$\square$] Créer backup complet base de données
    \item[$\square$] Tests d'intégrité référentielle
\end{itemize}

\clearpage

\subsection{MERCREDI 11 FÉVRIER - Jour 4 : Tests et Debugging}

\subsubsection{Matin (8h-12h) - Tests Intensifs}

\paragraph{Tous ensemble - Session Tests}
\begin{itemize}
    \item[$\square$] Chaque équipe teste fonctionnalités autres équipes
    \item[$\square$] Créer liste bugs : critiques / moyens / mineurs
    \item[$\square$] Tester avec questions réelles d'étudiants
\end{itemize}

\paragraph{Tests spécifiques}
\begin{itemize}
    \item[$\square$] \textbf{Backend :} Tester 100 questions variées, vérifier scores
    \item[$\square$] \textbf{Base de Données :} Test charge (100 requêtes simultanées)
    \item[$\square$] \textbf{Frontend :} Tester 3 navigateurs (Chrome, Firefox, Safari)
    \item[$\square$] \textbf{Structuration :} Vérifier qualité 1000+ Q/R finales
\end{itemize}

\subsubsection{Après-midi (14h-18h) - Corrections}

\begin{itemize}
    \item[$\square$] \textbf{Backend :} Corriger bugs critiques algorithme
    \item[$\square$] \textbf{Backend :} Améliorer gestion cas edge
    \item[$\square$] \textbf{Base de Données :} Optimiser requêtes lentes
    \item[$\square$] \textbf{Structuration :} Corriger Q/R problématiques
    \item[$\square$] \textbf{Frontend :} Corriger bugs UI/UX
    \item[$\square$] \textbf{Frontend :} Améliorer messages erreur
\end{itemize}

\subsubsection{Soir (18h-19h)}
\begin{itemize}
    \item[$\square$] Réunion priorisation : dernières tâches critiques
    \item[$\square$] Définir qui fait quoi pour jeudi/vendredi
\end{itemize}

\clearpage

\subsection{JEUDI 12 FÉVRIER - Jour 5 : Documentation et Démo}

\subsubsection{Matin (8h-12h) - Documentation Technique}

\paragraph{Équipe Backend}
\begin{itemize}
    \item[$\square$] Documenter API (README\_API.md) :
    \begin{itemize}
        \item Endpoints disponibles
        \item Format requêtes/réponses
        \item Exemples utilisation
    \end{itemize}
    \item[$\square$] Créer requirements.txt à jour
    \item[$\square$] Commenter code critique
\end{itemize}

\paragraph{Équipe Base de Données}
\begin{itemize}
    \item[$\square$] Documenter schéma BD (diagramme ER)
    \item[$\square$] Créer guide installation et migration
    \item[$\square$] Préparer scripts seed (données démo)
\end{itemize}

\paragraph{Équipe Structuration}
\begin{itemize}
    \item[$\square$] Rédiger guide contribution CSV
    \item[$\square$] Créer FAQ meta sur chatbot :
    \begin{itemize}
        \item "Comment fonctionne ce chatbot ?"
        \item "Que faire si pas de réponse ?"
    \end{itemize}
\end{itemize}

\paragraph{Équipe Frontend}
\begin{itemize}
    \item[$\square$] Créer page "À propos" expliquant projet
    \item[$\square$] Ajouter guide utilisation (tooltips, onboarding)
    \item[$\square$] Documenter code frontend
\end{itemize}

\subsubsection{Après-midi (14h-18h) - Préparation Démo}

\paragraph{Tous ensemble}
\begin{itemize}
    \item[$\square$] Créer scénario démonstration :
    \begin{enumerate}
        \item Introduction contexte
        \item Architecture technique
        \item Démonstration live (10 questions variées)
        \item Système feedback
        \item Statistiques
        \item Perspectives futures
    \end{enumerate}
    \item[$\square$] Préparer slides présentation (PowerPoint/Google Slides)
    \item[$\square$] Répéter démo 2-3 fois
\end{itemize}

\paragraph{Tâches spécifiques}
\begin{itemize}
    \item[$\square$] \textbf{Backend :} Préparer 10 questions démo impressionnantes
    \item[$\square$] \textbf{Base de Données :} Créer données feedback réalistes
    \item[$\square$] \textbf{Frontend :} Assurer visuel impeccable
    \item[$\square$] \textbf{Structuration :} Story de collecte données
\end{itemize}

\clearpage

\subsection{VENDREDI 13 FÉVRIER - Jour 6 : Finalisation et Livraison}

\subsubsection{Matin (8h-12h) - Polish Final}

\paragraph{Tous ensemble - Session Polissage}
\begin{itemize}
    \item[$\square$] Derniers ajustements UI/UX
    \item[$\square$] Vérification complète endpoints
    \item[$\square$] Test final end-to-end complet
    \item[$\square$] Vérification responsive (mobile, tablette, desktop)
\end{itemize}

\paragraph{Par équipe}
\begin{itemize}
    \item[$\square$] \textbf{Backend :} Logs et monitoring basique
    \item[$\square$] \textbf{Base de Données :} Backup final
    \item[$\square$] \textbf{Frontend :} Optimisation performances
    \item[$\square$] \textbf{Structuration :} Export CSV final
\end{itemize}

\subsubsection{Après-midi (14h-17h) - Déploiement \& Documentation}

\begin{itemize}
    \item[$\square$] Déployer sur serveur test (Heroku/PythonAnywhere/Railway)
    \item[$\square$] Finaliser README.md principal :
    \begin{itemize}
        \item Description projet
        \item Installation et setup
        \item Architecture
        \item Screenshots
        \item Guide utilisation
        \item Roadmap future
    \end{itemize}
    \item[$\square$] Créer CHANGELOG.md
    \item[$\square$] Tag version v0.1.0-prototype
\end{itemize}

\subsubsection{Fin de journée (17h-18h) - Démo Officielle}

\begin{important}[DÉMONSTRATION FINALE]
\begin{itemize}
    \item Présentation prototype devant équipe/encadrant/évaluateur
    \item Démonstration système complet
    \item Réponse aux questions
    \item Collecte feedbacks
    \item Discussion prochaines étapes
\end{itemize}
\end{important}

\clearpage

% Section 4 : Livrables
\section{Livrables Attendus (13 février 18h)}
\label{sec:livrables}

\subsection{Code}

\begin{itemize}
    \item[\faIcon{check}] Projet Django complet avec 3 apps (faq, chatbot, users)
    \item[\faIcon{check}] Frontend HTML/CSS/JS fonctionnel
    \item[\faIcon{check}] Base de données avec 1000+ Q/R
    \item[\faIcon{check}] API REST documentée
\end{itemize}

\subsection{Documentation}

\begin{itemize}
    \item[\faIcon{check}] README.md complet
    \item[\faIcon{check}] Documentation API (README\_API.md)
    \item[\faIcon{check}] Guide BD (schéma ER + migrations)
    \item[\faIcon{check}] Guide contribution CSV
    \item[\faIcon{check}] CHANGELOG.md
\end{itemize}

\subsection{Démo}

\begin{itemize}
    \item[\faIcon{check}] Application déployée et accessible en ligne
    \item[\faIcon{check}] Présentation PowerPoint (10-15 slides)
    \item[\faIcon{check}] Scénario démo préparé
    \item[\faIcon{check}] 10 questions test impressionnantes
\end{itemize}

\clearpage

% Section 5 : Indicateurs
\section{Indicateurs de Succès}
\label{sec:indicateurs}

\begin{table}[H]
\centering
\begin{tabular}{|p{4cm}|p{3cm}|p{5cm}|}
\hline
\rowcolor{blue!20}
\textbf{Critère} & \textbf{Objectif} & \textbf{Mesure} \\
\hline
Q/R dans la base & 1000+ & COUNT(*) FROM faq \\
\hline
Taux de réponse & >70\% & Questions score >0.6 \\
\hline
API fonctionnelle & 100\% & Tous endpoints testés \\
\hline
Interface utilisable & \faIcon{check} & Chat + feedback OK \\
\hline
Documentation & \faIcon{check} & README + API + BD \\
\hline
Démo prête & \faIcon{check} & Scénario testé 3x \\
\hline
\end{tabular}
\caption{Critères de validation du prototype}
\end{table}

\clearpage

% Section 6 : Points d'attention
\section{Points d'Attention Critiques}
\label{sec:attention}

\subsection{Risques identifiés}

\begin{enumerate}
    \item \textbf{Synchronisation équipes} : Réunions quotidiennes obligatoires (matin + soir)
    \item \textbf{Qualité vs Quantité} : Validation systématique 20\% des Q/R générées
    \item \textbf{Blocage technique} : Plan B si import CSV pose problème
    \item \textbf{Scope creep} : NE PAS ajouter fonctionnalités non prévues
    \item \textbf{Fatigue production} : Pauses régulières, rotation tâches
\end{enumerate}

\subsection{Bonnes pratiques}

\begin{itemize}
    \item[\faIcon{sync}] Commits Git fréquents (min. 2 par personne par jour)
    \item[\faIcon{comments}] Communication continue (groupe WhatsApp/Telegram)
    \item[\faIcon{bug}] Logger tous bugs dans fichier partagé
    \item[\faIcon{flask}] Tester après chaque feature importante
    \item[\faIcon{coffee}] Pauses régulières pour éviter fatigue
\end{itemize}

\subsection{Gestion Production Massive}

\begin{astuce}[Workflow Génération IA]
\textbf{Pour chaque vague de 100 Q/R :}

\begin{enumerate}
    \item Définir sous-thème précis
    \item Créer prompt IA optimisé
    \item Générer 120 Q/R (marge sécurité)
    \item Validation croisée : 20\% échantillon
    \item Correction anomalies détectées
    \item Export CSV consolidé
    \item Import dans base test
    \item Vérification vecteurs TF-IDF
\end{enumerate}

\textbf{Durée estimée par vague :} 1h30 - 2h
\end{astuce}

\clearpage


% Section 8 : Conclusion
\section*{Conclusion}

\vspace{1cm}

\begin{center}
\Large
\textbf{L'objectif est clair :}

\vspace{0.5cm}

\textbf{Un prototype fonctionnel et impressionnant} \\
\textbf{avec 1000+ Q/R en 6 jours}

\vspace{1cm}

\textbf{\faIcon{check} On peut le faire !} \\
\textbf{Let's build something amazing together !}

\vspace{2cm}

\rule{0.5\textwidth}{0.4pt}

\vfill

\normalsize
\textit{Document préparé par :} \\
\textbf{NKOUMOU TJADE} \\
Chef de la Division de Programmation \\
Club Informatique SUP'PTIC

\vspace{0.5cm}


\end{center}

\end{document}